\section{Related Work}
\paragraph{Evaluator reliability before adaptation.}
JudgeBench, RewardBench, and Plan-RewardBench expose failures of LLM judges and reward models on objectively or procedurally checkable inputs~\citep{tan2025judgebench,lambert2024rewardbench,wang2026planrewardbench}. They establish that a terminal signal can be wrong. Our question begins one step later: what happens when that signal chooses a persistent-state update and the resulting state is reused?

\paragraph{Feedback-conditioned memory construction.}
ReasoningBank distills structured reasoning memories from self-judged successful and failed experiences, while Reflexion stores evaluative verbal feedback across trials~\citep{ouyang2026reasoningbank,shinn2023reflexion}. Mem-$\alpha$ learns what to store from downstream reward, and AttriMem adds attribution-derived process feedback to address coarse outcome-level credit assignment~\citep{wang2025memalpha,li2026attrimem}. These works make feedback-conditioned construction established prior art. We instead freeze the construction substrate and source experience, use the terminal label as the sole treatment, and ask where that intervention remains distinguishable after storage.

\paragraph{Memory representation, retrieval, and reuse.}
Raw trajectories and long-context baselines remain important competitors for agent memory: ReasoningBank explicitly compares against raw experience, while ALMA and EvoMemBench evaluate learned or designed memories against trajectory-retrieval and strong long-context controls~\citep{ouyang2026reasoningbank,xiong2026alma,wang2026evomembench}. Very recent work further separates retrieval from post-retrieval reuse by holding the retrieved candidate and target state fixed while changing the support delivered to the policy~\citep{li2026beyondretrieval}. We therefore do not compare reward-conditioned memory only with omission, nor treat retrieval success as realized use. Our controls include the exact compact writer-input trajectory and an outcome-blind structured procedural rewrite using the same information and memory schema.

\paragraph{Reward-driven memory reuse and self-improvement variability.}
Memory Reward Inflation studies how proxy reward can overvalue stored memories and amplify reuse; RoMeRL targets misleading utility updates and the memory-reward trap; Ye et al. quantify variance, task-order sensitivity, and underspecification in self-improving agents~\citep{asadolahi2026inflation,yang2026romerl,ye2026fragility}. Their interventions operate at valuation, utility update, retrieval/reuse, or whole-system variation. Our treatment is earlier and more local: a controlled write-time label switch followed by a stage-resolved audit of exposure, policy uptake, and outcome.

\paragraph{Defended residual.}
The residual claim is therefore neither ``reward can be wrong'' nor ``feedback changes memory.'' It is that a byte-identical label intervention can be identified at the durable-state boundary, while downstream evidence must separately establish whether that difference is exposed and converted into branch-specific policy behavior. This separation lets positive forced-intervention evidence coexist with a negative native-transport boundary without contradiction.
